\documentclass[11pt]{article}
\usepackage[margin=0.9in]{geometry}
\usepackage{amsmath,amssymb,amsthm,booktabs}
\usepackage{mathrsfs}
\usepackage[T1]{fontenc}
\usepackage{lmodern,microtype}
\usepackage[colorlinks=true,linkcolor=blue,urlcolor=blue]{hyperref}
\newtheorem{theorem}{Theorem}[section]
\newtheorem{lemma}[theorem]{Lemma}
\newtheorem{proposition}[theorem]{Proposition}
\newtheorem{corollary}[theorem]{Corollary}
\theoremstyle{remark}
\newtheorem{remark}[theorem]{Remark}
\newcommand{\norm}[1]{\left\|#1\right\|}
\title{Toward a four-dimensional Yang--Mills construction\\
\large Notebook 16: local nonlinear responses and the classical coarse action}
\author{Research working document}
\date{16 October 2026}
\begin{document}
\maketitle
\begin{abstract}
Starting from the finite-regulator small-field minimizer of
Notebook 15, we prove exponential locality of its response to
coarse data at small nonconstant backgrounds. The bounds are
uniform in refinement and periodic volume, retain the same
physical H\"older exponent, and allow one weighted direction
at every analytic derivative order. The extra input is a local
physical-transpose derivative estimate for the nonlinear
constraint. An exact envelope identity then identifies the
derivative of the minimized classical action with the coarse
multiplier. This supplies a local nonlinear coarse Hessian,
its small perturbation from the flat Maxwell precision, and a
classical gauge identity. A separate exact Gaussian trace
calculation shows that a fixed physical smooth ball has vanishing
probability when the coupling decreases more slowly than the
square of the mesh spacing. This identifies a concrete need for
rescaled fluctuation domains. Notebook 17 supplies the resulting
fixed-step Gaussian domain and joint tails; Notebook 18 supplies
an overlapping-interaction integration theorem under explicit local
coefficient hypotheses. The complete model-specific integration
and iteration remain open.
\end{abstract}

\noindent\textbf{Status.}
Internally audited research proofs, with no external certification
or novelty claim. The prerequisite finite-regulator construction
and its ordinary-constraint transfer are given in Notebook 15.
The source provenance remains the deposited transcriptions at
\href{https://github.com/ethanwsussman/Yang--Mills-LLM-experiment/tree/e392ef45013a581664bab76b90867170b30fe10e}{commit
\texttt{e392ef45013a581664bab76b90867170b30fe10e}}.
No new original-paper verification is asserted here.

\section{The precise target}
\label{nr16-target}

Work with $SU(N_c)$ on compatible finite periodic lattices,
fine spacing $\eta=L^{-j}$ for fixed odd $L\ge3$, coarse
spacing one, and fixed $0<\alpha<1$. The background coarse
potential $B$ is arbitrary in a sufficiently small global
supremum ball; it need not be constant or decay at infinity.
The selected fine potential $A_*(B)$ uses the corrected
constraint and fixed flat gauge of Notebook 15. All derivatives
of this chosen representative below refer to that selection.

The ordinary-feasible representative from Notebook 15 is gauge
equivalent and has a separately controlled response through its
explicit gauge formula. Its leading extension generally differs
from $H$, while its minimized action value is exactly the same.
Keeping these coordinates distinct gives local response theorems
for both specified representatives and the ordinary-gauge
interpretation of the classical action.

The main analytic estimate puts an exponential weight on one
coarse direction while keeping the other directions and the
background unweighted. It yields decay away from a localized
insertion and uniform higher derivatives. A fixed-exponent
connected-tree bound involving all insertion sites at once is
a stronger statement and is not inferred from this estimate.

